% authority: science-draft
% status: draft/control
% route: refuter-stress-packet
% adoption_status: blocked_adoption_open_continuation
% refuter_result: scoped_obstruction
\documentclass[11pt]{article}
\usepackage[margin=1in]{geometry}
\usepackage[T1]{fontenc}
\usepackage{lmodern}
\usepackage{amsmath,amssymb,amsthm}
\usepackage{booktabs}
\usepackage{longtable}
\usepackage{xurl}
\usepackage[hidelinks]{hyperref}

\newtheorem{definition}{Definition}
\newtheorem{theorem}{Theorem}
\newtheorem{proposition}{Proposition}
\newtheorem{corollary}{Corollary}
\newtheorem{obstruction}{Scoped obstruction}
\newtheorem{nonconclusion}{Non-conclusion}
\newcommand{\Q}{\mathbb Q}
\newcommand{\Kspace}{\mathcal K_C}
\newcommand{\Qspace}{\mathcal Q_C}
\newcommand{\Aspace}{\mathcal A_C}
\newcommand{\Eop}{\mathcal E_C^{\alpha,\beta}}

\title{FiniteConstraintDynamicalViabilityStress\_v1\\
P8-T05 Constraint, Mode, Characteristic, and Stability Analysis}
\author{The \AE ther-Flow Research Project}
\date{2026-07-29}

\begin{document}
\sloppy
\maketitle

\section{Decisive result and authority boundary}

This packet stress-tests only the exact P8-T04 finite equation
\begin{equation}
 \Eop([h];u)=\alpha L_Ch-\beta L_Cu=0,
 \qquad
 [h]=\frac{\beta}{\alpha}[u],
 \label{eq:fixed}
\end{equation}
with the P8-T03 candidate restrictions \(\alpha>0\), \(\beta\in\Q\),
fixed source \(u\), and varied quotient slot
\(\Qspace=\Q^\Omega/\ker L_C\).  The result classification required by the
Refuter role is
\[
 \boxed{\texttt{scoped\_obstruction}}.
\]

The exact finite system has a complete algebraic constraint and mode count:
if \(n=|\Omega|\) and \(c=|\Pi_C|\) is the number of connected support
components, then it has \(c\) redundant component-shift directions,
\(n-c\) quotient coordinates, \(n-c\) independent algebraic equations, and
zero free homogeneous quotient-response modes after the fixed-source
constraint is imposed.  For \(\alpha>0\), its quotient Hessian is positive
definite.  This is a static convexity and solvability result.

The fixed record contains no time parameter, time derivative, kinetic form,
symplectic form, Hamiltonian constraint system, mass matrix, Lorentzian
background, causal cone, or nonlinear interaction.  Therefore physical
degrees of freedom, characteristic speeds, hyperbolicity, ghosts, tachyons,
gradient instabilities, and a strong-coupling scale are not defined by the
candidate.  A pair of explicit dynamical completions with the same static
equation but incompatible stability verdicts proves that no such verdict can
be inferred from Equation~\eqref{eq:fixed}.

The obstruction identifier is
\[
\texttt{OBST-P8T05-STATIC-FINITE-CONSTRAINT-DYNAMICS-UNDEFINED-001}.
\]
It is local to the unchanged P8-T04 candidate.  It is not a global no-go
theorem and does not establish future source-extension impossibility.

\section{Fixed finite system}

Let \(C=(P+P^{\mathsf T})/2\) be the symmetric conductance fixed by the
adopted P7 source package, and let \(L_C\) be its graph Laplacian.  If
\(\Pi_C\) is the set of connected components of the undirected support of
\(C\), define
\begin{align}
 \Kspace&=\ker L_C
 =\{k\in\Q^\Omega:k\text{ is constant on each }K\in\Pi_C\},\\
 \Qspace&=\Q^\Omega/\Kspace,\\
 \Aspace&=\Kspace^\perp.
\end{align}
The fixed-source action and its exact response operator are
\begin{align}
 S_C^{\alpha,\beta}([h];u)
 &=\frac{\alpha}{2}h^{\mathsf T}L_Ch
   -\beta(L_Cu)^{\mathsf T}h,\\
 \Eop([h];u)&=\alpha L_Ch-\beta L_Cu.
 \label{eq:action-response}
\end{align}
No target-side object enters this definition.  In particular, the row label
of the underlying finite source kernel is not target-spacetime time, and
\(L_C\) is not a spacetime differential or curvature operator.

\section{Algebraic constraints and mode count}

\begin{theorem}[Complete algebraic count]
\label{thm:count}
Let \(n=|\Omega|\), \(c=|\Pi_C|\), and \(r=n-c=\operatorname{rank}L_C\).
Then:
\begin{enumerate}
 \item the redundant representative \(h\in\Q^\Omega\) has \(n\)
 coordinates and exactly \(c\) component-shift directions in \(\Kspace\);
 \item the quotient carrier \(\Qspace\) has dimension \(r\);
 \item Equation~\eqref{eq:fixed} supplies exactly \(r\) independent
 algebraic equations in \(\Aspace\);
 \item the fixed-source solution set in the quotient is the singleton
 \([h]=(\beta/\alpha)[u]\); and
 \item the homogeneous linearized quotient equation has zero free solution
 dimension.
\end{enumerate}
\end{theorem}

\begin{proof}
The component indicators form a basis of \(\ker L_C\), so
\(\dim\Kspace=c\).  Rank--nullity gives
\(\dim\Qspace=\operatorname{rank}L_C=n-c=r\).  The restriction of \(L_C\)
to \(\Qspace\) is an isomorphism onto \(\Aspace\).  Hence
\(\alpha L_C[h]=\beta L_C[u]\) supplies \(r\) independent algebraic
equations and has the unique quotient solution in
Equation~\eqref{eq:fixed}.  Linearizing at that solution gives
\(\alpha L_C[\delta h]=0\), whose quotient kernel is trivial.
\end{proof}

\begin{corollary}[Spectral form]
\label{cor:spectral}
Choose an orthonormal quotient eigenbasis
\(\{v_a\}_{a=1}^{r}\) with \(L_Cv_a=\lambda_av_a\) and
\(\lambda_a>0\).  For
\[
 q:=h-\frac{\beta}{\alpha}u
 =\sum_{a=1}^{r}q_av_a,
\]
the exact quadratic displacement action and linearized equations are
\begin{align}
 S_C^{\alpha,\beta}([h];u)-S_C^{\alpha,\beta}
 \!\left(\frac{\beta}{\alpha}[u];u\right)
 &=\frac{\alpha}{2}\sum_{a=1}^{r}\lambda_aq_a^2,\\
 \alpha\lambda_aq_a&=0,\qquad a=1,\ldots,r.
 \label{eq:spectral}
\end{align}
For \(\alpha>0\), every quotient Hessian eigenvalue
\(\alpha\lambda_a\) is positive.
\end{corollary}

The phrases ``zero free homogeneous quotient-response modes'' and
``positive quotient Hessian'' are algebraic statements.  They do not say
that the model has zero propagating physical degrees of freedom.  A
propagating degree-of-freedom count requires a temporal phase space and a
constraint evolution law, neither of which is present.

\section{Exact two-state control}

For the P7 two-state forward control,
\[
 L_C=
 \begin{pmatrix}
 1/4&-1/4\\
 -1/4&1/4
 \end{pmatrix}.
\]
Its eigenvalues are \(0\) and \(1/2\).  The kernel is spanned by
\((1,1)^{\mathsf T}\), the quotient is one-dimensional, and the sole
quotient Hessian eigenvalue is \(\alpha/2>0\).  Writing
\(q=(z/2,-z/2)^{\mathsf T}\) gives
\[
 \frac{\alpha}{2}q^{\mathsf T}L_Cq=\frac{\alpha}{8}z^2,
 \qquad
 \frac{\partial}{\partial z}\frac{\alpha}{8}z^2
 =\frac{\alpha}{4}z.
\]
Thus the fixed-source solution is unique modulo the common shift.  No
frequency \(\omega\), wave number \(k\), characteristic speed, or kinetic
sign occurs in this exact calculation.

\section{Characteristic and hyperbolicity obstruction}

\begin{proposition}[No dynamical principal symbol is defined]
\label{prop:principal}
The linearized operator supplied by the fixed candidate is the finite
algebraic matrix
\[
 D_h\Eop=\alpha L_C.
\]
It is independent of a temporal frequency, spatial covector, background
metric, or differential order.  Consequently the candidate supplies no
principal symbol \(\sigma(\omega,\xi)\), no characteristic polynomial, no
characteristic cone, no propagation speed, and no hyperbolicity
classification.
\end{proposition}

\begin{proof}
A dynamical characteristic calculation requires the highest-derivative
part of an evolution equation and therefore a covector-dependent principal
symbol.  Equation~\eqref{eq:fixed} contains no derivative with respect to
any evolution parameter and no map from the finite support to a spacetime
cotangent bundle.  Calling the constant matrix \(\alpha L_C\) a
``principal symbol'' would not create the missing covector dependence or
evolution structure.
\end{proof}

The same missing datum blocks a Hamiltonian or gauge-invariant physical
mode count.  Component shifts are quotient redundancy of a static graph
functional; they have not been derived as first-class Hamiltonian
constraints or spacetime diffeomorphisms.

\section{Minimal dynamical-completion countermodels}

\begin{theorem}[Static-equation nonuniqueness]
\label{thm:countermodel}
Let \(M\) be any positive-definite matrix on \(\Qspace\), and let
\(q=h-(\beta/\alpha)u\).  The two task-local dynamical completions
\begin{align}
 \mathcal L_{+}(q,\dot q)
 &=\frac12\dot q^{\mathsf T}M\dot q
   -\frac{\alpha}{2}q^{\mathsf T}L_Cq,\\
 \mathcal L_{-}(q,\dot q)
 &=-\frac12\dot q^{\mathsf T}M\dot q
   -\frac{\alpha}{2}q^{\mathsf T}L_Cq
 \label{eq:completions}
\end{align}
have the same static stationary equation
\(\alpha L_Cq=0\), and hence the same fixed P8-T04 solution.  Their
dynamical verdicts are incompatible:
\begin{align}
 M\ddot q+\alpha L_Cq&=0
 &&\text{for }\mathcal L_{+},\\
 M\ddot q-\alpha L_Cq&=0
 &&\text{for }\mathcal L_{-}.
\end{align}
The first has positive kinetic energy and positive squared normal-mode
frequencies; the second has a negative kinetic form and exponentially
growing quotient branches.  Varying \(M\) changes the normal-mode
frequencies while preserving the same static equation.
\end{theorem}

\begin{proof}
The displayed Euler--Lagrange equations follow directly from
Equation~\eqref{eq:completions}.  Setting \(\dot q=\ddot q=0\) reduces both
to \(\alpha L_Cq=0\).  On the quotient, the eigenvalues of
\(\alpha M^{-1/2}L_CM^{-1/2}\) are strictly positive.  They are squared
frequencies for \(\mathcal L_{+}\) and growth-rate squares for
\(\mathcal L_{-}\).  The negative kinetic form of \(\mathcal L_{-}\)
also makes its Hamiltonian unbounded below.  Neither completion is adopted;
their sole use is to refute inference of a unique dynamical verdict from
the static equation.
\end{proof}

These countermodels add an artificial evolution parameter only for the
logical test.  They are not candidate repairs, source laws, physical
theories, or evidence that either branch describes the project.

\section{Stability phase diagram in the exact scope}

\begin{longtable}{p{0.18\linewidth}p{0.32\linewidth}p{0.40\linewidth}}
\toprule
Diagnostic & Exact finite result & Scope and non-conclusion\\
\midrule
Algebraic solvability &
Unique quotient solution for \(\alpha>0\) &
Exact static result; not a physical evolution theorem\\
Static Hessian &
\(\alpha L_C\) positive definite on \(\Qspace\) &
Convexity only; not a ghost or gradient-stability result\\
Redundancy &
\(c\) component shifts in representatives &
Not established first-class constraints or diffeomorphisms\\
Free algebraic response &
Zero-dimensional homogeneous quotient solution &
Propagating physical degree-of-freedom count is undefined\\
Principal symbol &
Absent &
Characteristic speeds and hyperbolicity are undefined\\
Kinetic matrix &
Absent &
Ghost status is undefined\\
Mass matrix and dispersion &
Absent &
Tachyon status is undefined\\
Spatial-gradient term on a spacetime &
Absent &
Gradient-instability status is undefined\\
Nonlinear interaction vertices &
Identically absent in the exact quadratic finite model &
No physical strong-coupling scale follows\\
\bottomrule
\end{longtable}

The coefficient and topology branches are:
\begin{enumerate}
 \item \(\alpha>0\), \(\beta\in\Q\): the admitted branch; the quotient
 Hessian is positive and the solution is unique.
 \item \(\beta=0\): the source decouples and the quotient solution is
 \([h]=0\); the dynamical classifications remain undefined.
 \item \(\alpha=0\): outside the admitted candidate; the response equation
 loses its invertible quotient Hessian and is generically inconsistent for
 nonzero \(\beta L_Cu\).
 \item \(\alpha<0\): outside the admitted candidate; the static quadratic
 form is unbounded below on every nontrivial quotient component.
 \item \(\lambda_\star\downarrow0\): a near-disconnection limit; the
 generic balanced susceptibility scales as
 \(1/(\alpha\lambda_\star)\), while the special source-generated forcing
 \(\beta L_Cu\) retains the exact response
 \([h]=(\beta/\alpha)[u]\).
 \item \(\lambda_\star=0\) after a topology change: the kernel and quotient
 dimension change.  The correct response is to rebuild the component
 quotient, not to call the jump a tachyon, ghost, or loss of
 hyperbolicity.
\end{enumerate}

\begin{proposition}[Near-disconnection control]
For the two-vertex Laplacian
\[
 L_\varepsilon=
 \begin{pmatrix}\varepsilon&-\varepsilon\\
 -\varepsilon&\varepsilon\end{pmatrix},
 \qquad \varepsilon>0,
\]
the nonzero eigenvalue is \(2\varepsilon\).  Under an arbitrary balanced
external forcing in that mode, the response norm scales as
\((2\alpha\varepsilon)^{-1}\).  Under the admitted source-generated
forcing \(\beta L_\varepsilon u\), the factor of \(\varepsilon\) cancels
and the exact solution remains
\([h]=(\beta/\alpha)[u]\).
\end{proposition}

The first behavior is algebraic ill-conditioning of a generic forced
system.  It is not a derived strong-coupling regime.  Strong coupling
requires a normalized kinetic sector and nonlinear interaction vertices,
both absent here.

\section{Scoped obstruction and freeze consequence}

\begin{obstruction}[Static finite constraint does not determine dynamical
viability]
\label{obs:dynamics}
The target claim under stress is: ``the exact P8-T04 equation determines a
dynamically viable gravitational sector.''  Its failed premise is the
absence of any source-derived temporal evolution law, kinetic form,
principal symbol, causal background, mass or gradient operator, and
interaction normalization.  The countermodel pair in
Theorem~\ref{thm:countermodel} shows that inequivalent healthy and
pathological dynamical completions share the same static equation.
Therefore the unchanged candidate does not determine physical degrees of
freedom, characteristic speeds, hyperbolicity, ghost status, tachyon
status, gradient stability, or strong coupling.
\end{obstruction}

This freezes the unchanged finite equation against reuse as evidence of
physical dynamical viability or a healthy GR limit.  It does not freeze
materially new source-side dynamics, a new kinetic or evolution law, a new
continuum construction, or independent audit.  P8-T06 may audit the exact
P8-T03--P8-T05 closure record after checkpoint, while every physical and
promotion block remains in force.

\begin{nonconclusion}
Positive static curvature is not ghost freedom.  A graph spectral gap is
not a mass gap, wave speed, or causal cone.  Component shifts are not
diffeomorphism gauge.  Algebraic near-singularity is not strong coupling.
The countermodel pair is not adopted physics.  No target metric,
stress-energy tensor, Einstein equation, exact-GR recovery, benchmark
promotion, canonical ontology edit, proof, publication, push, global no-go,
future source-extension impossibility, or completed derivation follows.
\end{nonconclusion}

\section{Internal source materials}

\begin{itemize}
 \item The exact P8-T04 finite equation, identity, specification, and
 controls.
 \item The P8-T03 finite candidate, coefficient ledger, and controls.
 \item The P8-T02 local effective-action closure target.
 \item The exact P7 finite source action and protected P7-T08 decision.
 \item The P6 Gate B separating certificate.
 \item The v21 P8-T05 plan and the active Distance-to-GR burden map.
\end{itemize}

No external source is used.  The dynamical completions are new
task-local countermodels derived in this packet.

\end{document}
